\documentclass[a4paper]{ltjsarticle}
\usepackage{amsmath,amssymb,mathtools}
\usepackage{geometry}
\geometry{top=12mm,bottom=12mm,left=15mm,right=15mm}
\usepackage{enumitem}
\setlist{nosep,leftmargin=1.6em,topsep=1pt}
\usepackage{titlesec}
\titleformat{\section}{\small\bfseries}{}{0em}{}[\vspace{-2pt}\hrule\vspace{1pt}]
\titlespacing{\section}{0pt}{5pt}{2pt}
\setlength{\parskip}{1pt}
\setlength{\parindent}{1em}
\pagestyle{empty}
\AtBeginDocument{\small}

\begin{document}

\begin{center}
  {\large\bfseries 屋内 RSSI 位置推定の精度改善に向けた手法提案 — 中間報告}\\[3pt]
  {\small 情報通信システム論Ⅰ \quad グループ [GROUP\_NAME] \quad [AUTHOR\_NAME] \quad 提出日：2026年7月3日}
\end{center}

\section{背景と目的}

豊橋技術科学大学の学内 WiFi（tutwifi / tutwifi2025）の RSSI のみを用いた屋内位置推定において，
ベースライン手法（WCL: 3.57 m，PBL: 4.38 m，CLA: 8.07 m \; 往路平均誤差）を \textbf{2 m 未満}に改善することを目標とする．
既存手法の根本的問題として，(1) 周波数帯混在による系統的バイアス，
(2) 選択 AP の凸包内に推定位置が拘束されること，
(3) WCL 重み関数の物理的根拠の薄さ，の 3 点を特定した．

\section{これまでにやったこと}

公表ベースライン PBL/CLA/WCL を生 RSSI から再計算し，icsR8\_text.txt Table 1 の Ave/Max/Std を
再現するライブラリ \texttt{icsr8} を実装した（Python，テスト付き）．
実装を通じて仕様書に未記載の暗黙の前処理（3F AP 限定，C 棟群のみ選択，物理 AP 単位集約，
tie-break 規則）を特定し，公表値 5 件の tie 事象（P19/P30/P35/P43/P49）を全て再現した．
さらに，改善手法の体系的調査ノートを整備し（Tier 1〜5，20 手法），本中間報告の提案に活用した．

\section{提案手法}

6/29 の対面ミーティングで確定した方針（廊下単一直線仮定・AP 数制限なし・統計的機械学習路線）
に基づき，以下の 3 手法を提案する．

\smallskip
\noindent\textbf{3.1 \; 周波数帯分離 + マルチバンド融合}

フリスの伝達公式より，2.4 GHz と 6 GHz の間には
$\Delta L = 20\log_{10}(f/f_0) \approx 8\,\mathrm{dB}$
の伝搬損失差がある（実測でも 2.4 GHz が 5 GHz より強く観測される傾向を確認）．
これを無視した従来手法は周波数混在による系統誤差を生む．
本手法では周波数帯ごとに経路損失パラメータ $P_{0,a,b},\,\alpha_{a,b},\,\sigma_{a,b}$ を独立推定し，
逆分散重み $w_b \propto 1/\hat{\sigma}_b^2$ でバンド融合する．
周波数ダイバーシティにより個別バンドのフェージング起因の分散を低減できる．

\smallskip
\noindent\textbf{3.2 \; 廊下座標系 1D fingerprinting（GP radio map + 階層推定）}

廊下は単一直線上にあるという仮定を活用し，2D 座標ではなく廊下 arc-length $s\in[0,S]$ 上で位置推定する．
まずセグメント分類（C 棟 / C2 棟 / C3 棟）を行い，次にセグメント内で 1D 回帰を行う階層構成とする．
1D では各 AP-band ごとの RSSI を Matérn-3/2 カーネルの Gaussian Process でモデル化する：
\[
  r_{a,b}(s) \sim \mathcal{GP}(m(s),\,k(s,s')), \quad
  k(s,s') = \sigma_f^2\!\left(1+\frac{\sqrt{3}|s-s'|}{\ell}\right)\exp\!\left(-\frac{\sqrt{3}|s-s'|}{\ell}\right).
\]
これにより 2 m 間隔の観測点間の連続補間と不確実性の定量化が可能となり，
角部（P16–P17 付近）でも C2 棟 AP の出現パターンによって識別できる．
Euclidean 距離でなく廊下 geodesic 距離を用いることで壁越しの疑似相関を回避する．

\smallskip
\noindent\textbf{3.3 \; 確率的 fingerprinting（重み最適化 + オフセット補正）}

RSSI の外れ値（人体遮蔽・フェージング）に対して Gaussian より頑健な Student-$t$ 分布を採用し，
AP 検出/非検出を二項分布でモデル化する（非検出を $-100\,\mathrm{dBm}$ 固定埋めで扱う WCL より統計的に正当）：
\[
  \log p(s\mid q) = \sum_{a,b}\beta_{a,b}\bigl[D\cdot\log t_\nu(r;\,\mu(s)+\hat{\delta}_q,\,\sigma(s)) + \log p(D\mid s)\bigr].
\]
共通オフセット $\hat{\delta}_q = \mathrm{median}_{a,b}[r_{q,a,b}-\mu_{a,b}(s)]$ を差し引くことで端末 AGC・
身体遮蔽・保持方向によるスキャン全体のシフトを吸収する（手法 4: Centered RSSI）．
加えて，10 回測定の分散が大きい AP-band 組に対し $w = w_{\rm base}/(1+\sigma^2/\sigma_{\rm ref}^2)$ で
下方重み付けを行い，安定 AP を優先する（手法 11: 分散ベース信頼度重み）．

\section{これからやること（スケジュール）}

\begin{itemize}
  \item W1（〜7/13）: 周波数帯分離パイプライン実装 + WCL アブレーション（帯域混在 vs 分離の定量比較）
  \item W2（〜7/20）: GP radio map / Student-$t$ fingerprinting 実装
  \item W3（〜7/24）: 3 手法のクロスバリデーション・結果まとめ・課題発表会スライド作成
\end{itemize}

評価プロトコル: 往路で DB 構築 → 復路でテスト（往復逆も実施）．
集計値（平均・分散・検出率）は train fold 内のみで計算し leakage を防止する．
報告指標は平均誤差・中央値誤差・90th パーセンタイル・2 m 以内率とする．

\section{現時点の課題と対策}

\begin{itemize}
  \item \textbf{小規模データの過学習}: 59 地点 × 10 回 × 48 AP では大規模 ML/RL は過学習リスクが高い
        → GP・Student-$t$ FP など説明可能な統計的機械学習で対応
  \item \textbf{L 字廊下の幾何的問題}: Euclidean 距離では壁越し点の疑似相関が生じる
        → corridor geodesic 距離（廊下グラフ上の最短路長）で代替
  \item \textbf{評価リーク防止}: 往路・復路の同一地点データを train/test に混ぜると「再認識」評価になる
        → 往復完全分離プロトコルで厳格に管理
\end{itemize}

\end{document}
